\chapter{Related Work and Positioning}

The mathematical treatment of symmetry belongs to the established tradition of equivariant bifurcation theory and invariant analysis rather than to an isolated geometry program. The group-theoretic organization of bifurcations, fixed-point subspaces, symmetry breaking, and mode interaction has been developed systematically in the equivariant bifurcation literature, especially by Golubitsky, Stewart, and Schaeffer \cite{golubitsky1988SingularitiesII}. The present work uses those general ideas for a specific descriptive question: what distinctions are available at a given invariant order, and how does the answer change when the invariant algebra is enriched? The $\ell=5$ quartic and sextic calculations should therefore be read as a concrete invariant-resolution result inside this broader mathematical tradition, not as a rediscovery of symmetry methods.

The oscillator examples belong to the classical synchronization literature initiated by Kuramoto \cite{kuramoto1975SelfEntrainment} and developed extensively in nonlinear dynamics. Their role here is deliberately narrow. They provide a transparent case in which collective motion can be separated from motion internal to the synchronized whole, thereby making the persistence-through-variation definition operational before the more complicated spherical calculations are introduced.

The philosophical framework intersects with structural realism, process philosophy, and dual-aspect traditions but is not reducible to any one of them. Plato supplies the historical vocabulary of intelligible form \cite{plato1997Complete}; Spinoza provides a canonical monist framework in which apparently distinct attributes belong to one substance \cite{spinoza1994Ethics}; and later process traditions emphasize becoming rather than static substance. These sources provide historical comparison rather than authority for the mathematical claims. The present framework treats relationalism as a methodological starting point whose ontology remains to be earned.

The consciousness program is positioned against several existing attempts to connect phenomenology and physical structure. Chalmers' formulation of the hard problem makes explicit the gap between functional explanation and phenomenal experience \cite{chalmers1995FacingUp}. Global neuronal workspace approaches formulate conscious access in terms of large-scale availability and ignition \cite{dehaene2011ConsciousProcessing}. Integrated information theory begins from proposed axioms of experience and derives requirements for physical cause--effect structure \cite{tononi2016IIT}. The present admissibility--occupancy distinction can be applied to all such theories: formal derivation establishes representational capability, while empirical occupancy requires independent evidence that the corresponding biological structure is necessary and sufficient for experience.

No-report paradigms attempt to separate neural activity associated with consciousness from activity associated with explicit report \cite{tsuchiya2015NoReport}. Subsequent critiques emphasize that removing overt report does not automatically remove attention, reflection, or other post-perceptual processes and may introduce different confounds \cite{duman2022NoReport}. The present framework therefore treats no-report methods as confound-ablation tools rather than privileged access to phenomenology.

Perturbational complexity offers a particularly useful empirical benchmark because it was designed to probe complexity associated with consciousness independently of ordinary sensory processing and behavioral report \cite{casali2013PCI}. It remains a physical measurement, but its perturbational character makes it valuable for the interventional stage of the evidential ladder. The proposed physical-sufficiency program uses such measures as components of progressively enriched physical descriptions rather than treating any single consciousness index as identical with experience.

Partial information decomposition originates in work by Williams and Beer on decomposing multivariate information into redundant, unique, and synergistic contributions \cite{williams2010PID}. The book uses PID as a method for testing informational complementarity, with the explicit warning that the result depends on the chosen redundancy measure and therefore requires estimator preregistration and robustness checks. Informational complementarity is not by itself metaphysical duality.

Finally, quantum cognition demonstrates that Hilbert-space probability models can capture contextual and order-dependent effects in cognition without requiring the claim that cognition is literally quantum physical \cite{pothos2022QuantumCognition}. This distinction motivates the quarantine of the antiunitary hypothesis in Appendix C. A complex projective model of cognition may be useful descriptively; an antiunitary physical--experiential bridge requires a much stronger demonstration that genuinely complex phenomenological geometry is independently necessary.

Across these literatures the book's intended contribution is not a new universal mathematics that replaces existing theories. It is a cross-cutting discipline for keeping several inferential steps separate: the expressive capacity of a mathematical language, the empirical occupation of structures admitted by that language, and the ontological interpretation of a successful model. The three governing principles are designed as reusable constraints on theories that otherwise move too easily from formal elegance to claims about reality.
